% 本文件由 LaTeX 排版 Agent 根据有效 translation.json 自动生成。
% author=Gregory of Nyssa; work_id=2913; title=论朝圣

\chapter{论朝圣}

% structure: p0001 source=h1 target=omitted reason=page-title-already-rendered-by-parent

由于，我的朋友，你在信中向我提出了一个问题，我认为我有责任按恰当的顺序回答你关于这问题的一切相关点。那么，我的意见是，对于那些已一次而永远献身于更高生活的人来说，持续关注福音中的言论是件好事，并如同那些用规矩校正任何给定材料中工作的人，藉着规矩的笔直，将他们双手发现的弯曲矫正为笔直，同样，我们应当对这问题应用一种严格无瑕的尺度——我指的是福音的生活准则——并按照它，在天主面前引导我们自己。如今，在那些已进入隐修和独修生活的人中，有些人将敬拜耶路撒冷那些能见到我主在世生活的纪念地视为他们虔诚的一部分；那么，最好查看这准则，若其诫命指示遵守这些事，便当做主的实际命令来履行这工作；但若它们完全超出师傅的诫命，我看不出有什么命令能使一个已成为义务之律法的人去热心地履行其中任何一件。当主邀请蒙福者进入天国继承时，祂并未将耶路撒冷朝圣列入他们的善行中；当祂宣布真福时，祂并未将那种虔诚列入其中。但对于那既不使我们有福、也不使我们进入天国道路的事，为何要追求它，让智者思考吧。即使他们所做的有些益处，但即便如此，成全的人最好也不要急于实践它；然而，当仔细审视此事时，发现它对那些已开始更严格生活的人造成道德上的伤害，它远不值得认真追求，实际上需要极大的谨慎，以防止那已献身于天主的人受到它任何有害影响。那么，它有害之处是什么呢？圣洁生活向所有人开放，无论男女。那默观生活的独特标志是端庄。但端庄保持在分开而独立生活的团体中，这样就不应有男女相遇和混合；男人不应在妇女陪伴下匆忙遵守端庄的规则，妇女也不应在男人陪伴下如此。但旅途的必需品常常容易使这种谨慎在遵守这些规则方面变得非常随便。例如，一个女人不可能在没有向导的情况下完成如此长途的旅行；由于她天然的软弱，她必须被扶上马又被扶下；在困难的地方她需要被支撑。无论我们设想她有一个熟人做这种忠心的服务，还是雇一个侍从来完成，两种方式都不能免于受责备；无论她依靠陌生人的帮助，还是依靠自己的仆人，她都无法遵守正确行为的法则；而且，由于东方的旅店、客栈和城市提供了许多放纵和对恶行漠不关心的例子，一个人穿过这样的烟尘，怎么可能不眼睛刺痛呢？当耳朵和眼睛被玷污，心也被沾染，通过眼和耳接受了所有那些污秽，怎么可能穿过这些传染之地而不受感染呢？此外，到达那些著名地点本身的人能得到什么益处呢？他不能想象我主现在仍以身体生活在那里，而是已离开我们这些外邦人；也不能想象圣神在耶路撒冷丰盛，却无法到达我们这里。然而，若真能通过可见的标记推断天主的存在，那么更有理由认为祂居住在卡帕多细亚民族而非任何境外之地。因为那里有多少祭台，我主的名在其中受光荣！全世界的其余地方也难以数出这么多。再者，若神圣恩宠在耶路撒冷附近比别处更丰盛，那么那里的居民就不会如此盛行罪恶；但事实上，在他们中没有任何一种不洁不被实行；奸诈、奸淫、偷盗、偶像崇拜、下毒、争吵、凶杀盛行；最后一种邪恶如此普遍，以至于世界任何地方的人都不像那里那样随时准备互相残杀；亲属像野兽一样互相攻击，为了无生命的掠夺而流血。那么，在发生这种事情的地方，请问，你从哪里证明神圣恩宠的丰盛呢？但我知道许多人会反驳我所说的一切；他们会说：\Quote{为什么你自己不也为自己定下这准则呢？若敬虔的朝圣者因去过那里而无所获益，你为何要承受如此长途跋涉的辛劳呢？}让他们听我的辩解。由于我生命支配者使我担任的职务所需，为了圣议会所决定的纠正，我有责任访问阿拉伯教会所在之处；其次，由于阿拉伯毗邻耶路撒冷地区，我承诺也与耶路撒冷圣教会的首领们商谈，因为他们的事务混乱，需要一位仲裁者；第三，我们最虔敬的皇帝通过驿马提供了旅途便利，以致我们无需忍受我们在别人身上注意到的那些不便；实际上，我们的马车对我们来说就像一座教堂或隐修院，因为我们在整个旅程中都歌唱圣咏并在主内守斋。因此，让我们自己的情况不使任何人困惑；反而要更多地听从我们的劝告，因为我们是就自己亲眼所见之事给出建议。我们宣认那显现的基督是真天主，无论是在耶路撒冷停留之前还是之后；我们对祂的信德既未增加也未减少。在看到白冷之前，我们就已知道祂通过童贞女降生成人；在看到祂的坟墓之前，我们就已相信祂从死者中复活；未看到橄榄山，我们就宣认祂升天是真实的。我们从那里的旅行只得到这一点益处，即通过比较，我们认识到我们自己所在的地方远比境外之地更神圣。因此，敬畏主的人啊，在你们所在之处赞美祂。地点的改变并不能使人更接近天主，无论你在哪里，天主都会来到你身边，只要你灵魂的内室被发现适合祂居住并在你内行走。但若你使你的内人充满邪恶的思想，即使你在哥耳哥达，即使你在橄榄山上，即使你站在复活纪念石上，你也将如一个尚未开始宣认祂的人一样，远离接受基督进入你内。因此，我亲爱的朋友，劝告弟兄们离开身体去往主那里，而不是离开卡帕多细亚去往巴勒斯坦；若有人引用我主对门徒所说的不可离开耶路撒冷的命令，就让他明白其真正含义。因为圣神的恩赐和分施尚未降临到宗徒们身上，我主命令他们留在原地，直到领受从上而来的能力。现在，若起初当圣神以火焰形像分施各恩赐时发生的事一直持续至今，那么所有人都应留在那分施发生的地方；但若圣神\BibleQuote{风随意吹}，那么在这里成为信徒的人也分享那恩赐；且是按他们信德的比例，而非由于他们去耶路撒冷朝圣。\par

\SourceRecord{Translated by William Moore and Henry Austin Wilson.；Nicene and Post-Nicene Fathers, Second Series；Vol. 5.；Edited by Philip Schaff and Henry Wace.；Buffalo, NY: Christian Literature Publishing Co.,；1893.；<http://www.newadvent.org/fathers/2913.htm>.}
